% ---------------------------------------------------------------------
%  Chapitre 6 — Conclusion : bilan + ouverture
%  (les limites détaillées et les perspectives vivent au chap. 5 ;
%   ici, on répond à la question du chap. 1 et on referme la boucle)
% ---------------------------------------------------------------------
\chapter{Conclusion}
\label{chap:conclusion}

L'objectif fixé au premier chapitre n'était pas de découvrir une stratégie
gagnante, mais de construire un protocole capable d'en juger une quelconque.
Ce protocole est désormais construit, éprouvé, et il a rendu son verdict~:
aucune des dix stratégies d'analyse technique ne démontre d'avantage réel une
fois ses conditions d'application vérifiées.

Le résultat central mérite d'être rappelé. Le test de référence déclarait deux
stratégies gagnantes. La seule vérification de la condition d'indépendance a
renversé les deux verdicts~: l'une se révélait un faux positif né de la
répétition d'une même information sur des milliers de trades, l'autre un
avantage réel mais trop épisodique dans le temps pour être retenu. Dans le
même temps, le témoin construit pour être gagnant a franchi chaque étage sans
faiblir. Un outil de validation se juge à sa capacité d'attraper ses propres
erreurs tout en reconnaissant le vrai~: la démonstration est faite.

Une idée traverse l'ensemble de ce travail. La dépendance des données y a joué
deux rôles opposés. Simultanée, portée par le facteur marché qui explique près
des trois quarts des mouvements, elle était la nuisance qui faussait les
tests, et il a fallu la mesurer puis la neutraliser pour juger proprement.
Retardée, elle aurait au contraire été la seule ressource qu'une stratégie
puisse exploiter, et le marché n'en contient presque pas, conformément à
l'efficience faible. Neutraliser l'une, chercher l'autre~: c'est la même
statistique, vue des deux côtés du même objet.

Le protocole est ainsi prêt à juger les prochains signaux de l'agent, en
commençant par leur validation temporelle.
